\scanpage{121}
\mspage{60}
Ceci posé, $\mathcal E_U^{\overline L}$ se récupère à partir de $\mathcal E_L^{\overline L}$, comme quotient de ce dernier, en prenant \uline{tous} les $V$ de $L$ ayant un centre sur $U$ (il suffit \unsure{=|même} de prendre les $V=\underline{O}_{U,x}$, où $x$ est un pt de codim 1 dans $U$), et en divisant par le ss-groupe fermé invariant engendré par les $I_{x'}$ correspondants.

On peut regarder l'espace loc. annelé, $\varprojlim$ des modèles \inserted{(réguliers)} propres de $L$ sur $K$
\[
(31)\qquad\mathfrak X_{L/K}\simeq\varprojlim_{\substack{X\text{ modèle}\\\text{propre de }L/K}}X
\]
\margin{\uline{NB} Cette définition ne dépend pas de l'hypothèse de résolution des sing --- car on peut prendre des modèles normaux --- pas néc. réguliers --- (Ça revient au \unsure{=|même}, si on dispose de la résolution\ldots)}
--- ce n'est \uline{pas} un schéma ! \inserted{(si $n\geq2$)}. \deleted{\unsure{L'on}} \inserted{Si}
$x\in\mathfrak X$, \deleted{\unsure{l'anneau de}} \inserted{$k(x)$} est une ext. de $K$ de degré de transcendance $d\leq n$ --- on l'appelle \uline{dimension} de $x$, \inserted{et $n-d$ est la \uline{codim.}} Il y a un pt et un seul de codim 0 --- correspondant aux pts gén. des $X$. Les points de codim 1
\margin{120}
